\documentclass[11pt,a4paper,sans]{moderncv}

% --- CONFIGURACIÓN DE COLORES SOBRIOS ---
\moderncvstyle{classic} 
\moderncvcolor{grey} % Base gris

% Aquí personalizamos el color exacto de las líneas y títulos:
\definecolor{color1}{RGB}{60, 60, 60} % Gris Oscuro Grafito
% Si prefieres Azul Marino, comenta la de arriba y usa esta:
% \definecolor{color1}{RGB}{0, 32, 96} 

% --- PAQUETES ---
\usepackage[T1]{fontenc}
\usepackage[utf8]{inputenc}
\usepackage[spanish,es-noshorthands]{babel}
\usepackage[scale=0.82]{geometry} 

% --- DATOS PERSONALES (Tal cual tu PDF) ---
\firstname{IGNACIO}
\familyname{VALMASEDA}
\title{Especialista en Riesgos y Diseño Metodológico}
\address{Las Palmas / Madrid}{España}

\extrainfo{
    \emailsymbol\ \href{mailto:ignacio.valmaseda@protonmail.com}{ignacio.valmaseda@protonmail.com}\\
    \mobilesymbol\ (+34) 664 645 676
}

\begin{document}

\makecvtitle

\section{PERFIL PROFESIONAL}
\cvitem{}{Especialista en el diseño metodológico de operaciones financieras, definiendo los estándares de cálculo de rentabilidades, garantías y parámetros de producto. Experto en modelización compleja para la gestión de capital y valoración de activos mediante simulaciones y contraste de modelos. Perfil con gran capacidad analítica, enfocado en la precisión técnica, el desarrollo de informes de control y la documentación metodológica detallada.}

\section{EXPERIENCIA PROFESIONAL}

\cventry{2025--Actualidad}{Especialista de Riesgos}{Banco Santander}{Madrid}{}{
\begin{itemize}
    \item \textbf{Diseño Metodológico y Soporte Experto:} Definición matemática de los algoritmos de cálculo (rentabilidad, garantías, amortización) y resolución de consultas complejas sobre los motores de cálculo. Referente para el desarrollo de mejoras metodológicas y la interpretación técnica de resultados en la operativa del banco.
    \item \textbf{Valoración y Herramientas Cuantitativas:} Dominio de herramientas de valoración avanzada, modelos de VaR y matrices de correlación de activos; capacidades críticas para la gestión y calibración de modelos de Capital y Pensiones mediante simulaciones de Monte Carlo.
    \item \textbf{Optimización de Algoritmos:} Rediseño y simplificación de ecuaciones de calibración y procesos de cálculo, logrando optimizar la eficiencia técnica y reducir tiempos de ejecución de varias semanas a 24--48 horas.
    \item \textbf{Documentación Técnica y Regulación:} Elaboración de documentación técnica y matemática de modelos, informes de backtesting y análisis ad-hoc, asegurando la convergencia con el marco regulatorio vigente.
\end{itemize}}

\vspace{2mm}

\cventry{2018--2025}{Analista de Riesgos}{Banco Santander}{Madrid}{}{
\begin{itemize}
    \item \textbf{Soporte Metodológico:} Soporte técnico en la definición de algoritmos de cálculo y lógica financiera de las operaciones.
    \item \textbf{Optimización y Eficiencia:} Rediseño de procesos y mejora de algoritmos, centrando esfuerzos en la automatización de herramientas críticas.
    \item \textbf{Análisis Regulatorio:} Documentación técnica y matemática con enfoque en la aplicación de la regulación bancaria vigente.
    \item \textbf{Desarrollo Cuantitativo:} Creación de aplicaciones en R/Shiny para monitorización de KPIs y soporte en análisis de Capital.
\end{itemize}}

\vspace{2mm}

\cventry{2018}{Becario en Análisis Cuantitativo}{Banco Santander}{Madrid}{}{
\begin{itemize}
    \item \textbf{Formación Técnica:} Formación inicial en el funcionamiento de los motores de cálculo, la lógica de riesgos y la estructura de bases de datos del banco.
    \item \textbf{Soporte Técnico:} Colaboración en la validación de herramientas internas y en la automatización de reportes de seguimiento y carga de datos.
    \item \textbf{Finalización de Beca:} Incorporación a la plantilla fija del departamento tras cuatro meses de beca.
\end{itemize}}

\section{FORMACIÓN ACADÉMICA}
\cventry{2017--2018}{Máster en Banca y Mercados Financieros}{U. Nebrija / C.E. Garrigues}{Madrid}{}{Banca de Inversión, valoración de activos, derivados, gestión de carteras y análisis de riesgos.}
\cventry{2016--2017}{Máster en Bioestadística}{U. Complutense de Madrid}{Madrid}{}{Modelización estadística avanzada, procesos estocásticos, análisis multivariante y programación en R.}
\cventry{2010--2016}{Grado en Física}{U. Complutense de Madrid}{Madrid}{}{Especialidad en Física Aplicada (Geofísica y Meteorología): análisis de sistemas complejos y modelos físico-matemáticos.}

\section{HABILIDADES TÉCNICAS}
\cvitem{Definición Metodológica}{Diseño y definición de lógica de operaciones, estándares de cálculo de rentabilidad, estructuras de amortización y parámetros técnicos de producto.}
\cvitem{Modelización y Análisis}{Modelización matemática compleja, simulación de Monte Carlo, procesos estocásticos y validación de modelos (backtesting).}
\cvitem{Software y Programación}{R, Shiny, SQL, Excel/VBA avanzado y Python.}
\cvitem{Idiomas}{Español (Nativo), Inglés (C1 - Certificación Cambridge BEC Higher).}

\end{document}